\section{An intrinsic coefficient principle}
\label{sec:structural-coefficient-principle}

The web problem combines two operations: extracting coefficient data
from an abstract scheme and reconstructing a relation space from those
data.  The first operation admits a formulation independent of webs or
quadratic multiplication.  We state the criterion and retain its
polar-restriction application for linear systems in every dimension.
Theorem~\ref{thm:automatic-quadratic-coefficients} subsequently derives
all of its geometric hypotheses directly for a full Grassmannian
family of the original quadratic multiplication schemes.

\subsection{The structural criterion}

Let \(X\) be a finite-type complex scheme with a closed subvariety
\(B\) specified by a construction preserved by every abstract scheme
isomorphism.  Write \(\mathfrak b\subset\OO_X\) for its ideal and
\[
 \mathcal G=\bigoplus_{r\ge0}\mathfrak b^r/\mathfrak b^{r+1},
 \qquad \mathcal A=\Sym_{\OO_B}(\mathcal G_1),
 \qquad \mathcal I=\ker(\mathcal A\twoheadrightarrow\mathcal G).
\]
Assume the following conditions, for fixed integers \(n\ge2,m\ge1\).

\smallskip
\noindent\textup{(i)} \(B\) is connected and projective,
\(H^0(B,\OO_B)=\C\), and \(\mathcal G_1\) is locally free.
Its dual admits a presentation \(V\otimes\mathcal U^*\), where
\(\dim V=\rank\mathcal U=n\).

\smallskip
\noindent\textup{(ii)} The reduction of \(\Spec_B\mathcal G\)
in its degree-one ambient bundle is the generic determinant cone.
The bundle \(\Pj_B(\mathcal U^*)\) is not projectively trivial.
Thus precisely one of the two ruling parameter bundles of its rank-one
locus is projectively trivial.

\smallskip
\noindent\textup{(iii)} If \(\mathcal D\subset\mathcal A\)
is the ideal of that reduced cone and \(\mathcal L=\mathcal D_n\),
then \(\mathcal I=\mathcal D\mathcal J\), where
\(\mathcal J\) is generated in degree \(m\).  Under the Cauchy
decomposition in any local frame of \(\mathcal U\),
\begin{equation}\label{eq:structural-coefficient-hypothesis}
 \mathcal J_m=
 \bigoplus_{\lambda\in\Lambda}
          L_\lambda\otimes\mathbb S_\lambda\mathcal U
 \ \subset\
 \Sym^m(V^*\otimes\mathcal U),
 \qquad L_\lambda\subset\mathbb S_\lambda V^*,
\end{equation}
where \(\Lambda\) is a specified collection of partitions of \(m\)
of length at most \(n\), and the left subspaces are constant over
\(B\).  Zero subspaces are permitted.  The Cauchy inclusions, not
just their representation labels, are the coefficient realizations
\(\beta\otimes\xi\mapsto[T\mapsto\beta(\mathbb S_\lambda(T)\xi)]\).

These conditions are hypotheses on the normal cone and its residual
coefficient module.  They do not assume that an abstract isomorphism
extends to an ambient space, preserves a chosen frame, or already acts
on all the left subspaces through a common projective transformation.

\begin{theorem}[Intrinsic coefficient reconstruction]
\label{thm:structural-coefficient-reconstruction}
Suppose \(X,B\) satisfy \textup{(i)--(iii)}.  The abstract scheme
\(X\) determines the collection \((L_\lambda)_{\lambda\in\Lambda}\)
up to one common projective coordinate transformation on \(V\).
More precisely, an isomorphism between two such schemes induces a
constant projective isomorphism \([g]:\Pj(V)\to\Pj(V')\) such that
\[
                L_\lambda=\mathbb S_\lambda(g)^*L'_\lambda
                       \quad(\lambda\in\Lambda).
\]
The equality means equality of subspaces; scalar changes of a linear
lift of \([g]\) do not affect it.
\end{theorem}
\begin{proof}
An isomorphism transports \(B\), every power of \(\mathfrak b\),
\(\mathcal G_1\), and the surjection
\(\mathcal A\twoheadrightarrow\mathcal G\).  In particular it
transports \(\mathcal I\) in the intrinsic degree-one ambient
bundle.  Iterated reduced singular loci of the determinant cone give
its rank-one locus.  Lemma~\ref{lem:rank-one-fano-scheme} reconstructs
its two ruling parameter schemes, including their scheme structure.
Condition \textup{(ii)} distinguishes the trivial ruling intrinsically.
Lemma~\ref{lem:oriented-segre-descent} then gives one constant
projective map \([g]\).  After choosing a constant linear lift and
local frames on both bases, the ambient bundle map has the form
\(T\mapsto gTh^{-1}\).  These local choices are made only after
constructing the graded algebra, its degree-one part, and the ruling.

Locally \(\mathcal D=(d)\) with \(d=\det T\).  Its leading
coefficient is a unit, so it is a nonzerodivisor over every coefficient
ring.  Hence \((\mathcal I:\mathcal D)=\mathcal J\), and
\[
 \mathcal J_m\simeq\mathcal L^*\otimes\mathcal I_{n+m}.
\]
This recovers the residual module without choosing a scalar determinant
equation.  Under \(T\mapsto gTh^{-1}\), the Cauchy coefficient
map transports a tensor by
\(\mathbb S_\lambda(g)^*\otimes\mathbb S_\lambda(h^{-1})\).
Projection to the specified Cauchy summand, followed by contraction
against its whole right dual, therefore transports exactly
\(L'_\lambda\) to \(L_\lambda\).  The right map is invertible
and has no effect on that left image.  Both determinant-line transition
functions and changes of a lift of \([g]\) act by scalar units.
Since every partition has size \(m\), this ambiguity is common
in degree and disappears on subspaces.  This proves the assertion.
\end{proof}

\begin{corollary}[The web instance of the criterion]
\label{cor:structural-web-instance}
For \(R\in G_4^\circ\), the scheme \(\widehat D_R\) satisfies
\textup{(i)--(iii)} with
\[
 B=\Gr(4,S_R),\qquad n=4,\qquad m=12,\qquad
 \Lambda=\{(5,4,2,1),(4,4,4,0)\}.
\]
The left subspaces are the two lines
\(\C\beta_{R,\lambda}\).  Their intrinsic recovery, exterior
duality, and the support gap imply Theorem~\ref{thm:main-web-reconstruction}.
\end{corollary}
\begin{proof}
The intrinsic base and determinant normal cone are
Lemma~\ref{lem:deepest-normal-cone}.  The nontriviality of the other
ruling is its Schubert-line test.  The residual factorization is
Lemma~\ref{lem:intrinsic-residual-ideal}, and
\eqref{eq:structural-coefficient-hypothesis} is exactly
Lemma~\ref{lem:universal-schur-readout}.  Exterior duality is
Lemma~\ref{lem:readout-exterior-twist}.  The final uniqueness is
Theorem~\ref{thm:global-schur-recombination}; the smooth Jacobian
has four essential variables.  Thus every hypothesis is verified on
the full stated smooth locus, not on a newly restricted generic subset.
\end{proof}

\subsection{Determinantal failure for polar restrictions}

Fix \(n\ge2\), \(d\ge1\), a vector space \(H\) of dimension
\(s>n\), and an \(n\)-dimensional vector space \(V\).  Put
\(B=\Gr(n,H)\), with tautological bundle \(\mathcal U\), and
\(M=\Tot\Hom(\mathcal U,V)\).  Let
\(0\ne L\subset\Sym^dV^*\) be a linear system, not furnished with
a preferred basis.  The tautological map \(T:\mathcal U\to V\)
on \(M\) gives the polar-restriction evaluation
\[
 \operatorname{ev}_{L,T}:L\otimes\Sym^d\mathcal U\longrightarrow\OO_M,
 \qquad f\otimes\xi\longmapsto f(\Sym^d(T)\xi).
\]
Here and below bundles from \(B\) are pulled back to \(M\).
Define a vector-bundle map and its zeroth-Fitting failure scheme by
\begin{align}
 \Psi_L&=\begin{pmatrix}\det T&0\\0&\operatorname{ev}_{L,T}\end{pmatrix}:
   \det\mathcal U\oplus(L\otimes\Sym^d\mathcal U)
          \longrightarrow\det V\oplus\OO_M,
             \label{eq:polar-fitting-map}\\
 Z_L&=V(\Fitt_0\coker\Psi_L)\subset M.
             \label{eq:polar-fitting-scheme}
\end{align}
In local frames, if \(f_1(T),\ldots,f_a(T)\) are the entries
of the evaluation row and \(\delta=\det T\), this is literally
\[
 \Psi_L=\begin{pmatrix}
    \delta&0&\cdots&0\\
    0&f_1(T)&\cdots&f_a(T)
 \end{pmatrix}.
\]
Its only nonzero maximal minors are \(\delta f_i(T)\).
The two blocks ask simultaneously whether \(T\) is invertible and
whether the restricted system is nonzero.  This is a different
failure map from cube-zero multiplication; neither a relation web nor
its Jacobian covariant is used in its definition.

\begin{theorem}[Reconstruction of linear systems from polar failure]
\label{thm:polar-system-torelli}
For fixed \(n,d,s\) as above and any nonzero systems
\(L\subset\Sym^dV^*\), \(L'\subset\Sym^dV'^*\),
\begin{equation}\label{eq:polar-system-torelli}
 Z_L\simeq Z_{L'}\quad\Longleftrightarrow\quad
           g^*L'=L\text{ for a linear isomorphism }g:V\to V'.
\end{equation}
The isomorphism on the left is abstract: it need not preserve the
projection to \(B\), the zero section, the grading, or an ambient
embedding.  All reductions \((Z_L)_{\mathrm{red}}\) are the same
determinant hypersurface in \(M\).  Thus the nonreduced scheme,
not its reduction, separates the projective-equivalence classes of
these linear systems.  No smoothness hypothesis on the hypersurfaces
or their base locus is required.
\end{theorem}
\begin{proof}
In a frame, write \(J_L\) for the ideal generated by all
coefficients of \(f\circ T\), for \(f\in L\).  The maximal
minors of the block map \eqref{eq:polar-fitting-map} give
\begin{equation}\label{eq:polar-residual-factorization}
                 I(Z_L)=(\det T)J_L.
\end{equation}
No zero of \(J_L\) has invertible \(T\): an invertible
substitution cannot annihilate a nonzero polynomial.  Therefore
\(V(J_L)\) is contained in the determinant hypersurface, and the
reduction in \eqref{eq:polar-residual-factorization} is precisely
that hypersurface with its reduced structure.  This also follows
locally from the primality of the generic determinant and the
Nullstellensatz, since \(B\) is smooth.

The reduced singular locus of the rank-at-most-\(r\) matrix
variety, for \(0<r<n\), is the rank-at-most-\(r-1\) variety.
For clarity, invert an \(r\)-minor and use block elimination:
the remaining lower-right block is zero, giving a smooth chart.
At smaller rank a similar block reduction leaves minors of size
at least two, with zero linear terms and tangent dimension exceeding
the local dimension.  These charts prove the assertion, also after
product with a smooth base.  Applying it repeatedly to
\((Z_L)_{\mathrm{red}}\) recovers its zero section \(B\)
intrinsically.  Its normal-cone algebra is the original graded algebra
\(\Sym(V^*\otimes\mathcal U)/((\det T)J_L)\): the ideal is
homogeneous and generated in degree \(n+d\).  In particular its
degree-one conormal is \(V^*\otimes\mathcal U\), before any
factorization is chosen by an isomorphism.

The two ruling parameter bundles are \(\Pj(V)\times B\) and
\(\Pj_B(\mathcal U^*)\).  Only the first is projectively trivial.
To check the second, restrict to a Schubert line in \(\Gr(n,H)\).
There \(\mathcal U^*=\OO^{n-1}\oplus\OO(1)\).  If its
projectivization were trivial, the vector bundle would be a direct
sum of \(n\) copies of one line bundle: pull back the relative
\(\OO(1)\) and push it down to see this assertion.  Its determinant
would have degree divisible by \(n\), whereas the displayed
splitting has determinant degree one.  This is impossible for
\(n\ge2\).

Finally the homogeneous coefficient map
\begin{equation}\label{eq:polar-cauchy-inclusion}
 \Sym^d V^*\otimes\Sym^d\mathcal U
       \longrightarrow\Sym^d(V^*\otimes\mathcal U)
\end{equation}
is the \((d)\)-Cauchy inclusion.  It is injective even before the
Cauchy decomposition is invoked.  Evaluate on a rank-one map
\(T=v\otimes a\), with \(a\in\mathcal U^*\): a tensor
\(f\otimes\xi\) evaluates to \(f(v)a^d(\xi)\).  These products
identify the tensor product with the polynomials bihomogeneous of
degree \((d,d)\) in \((v,a)\), so a zero coefficient polynomial
has zero tensor.  Consequently
\begin{equation}\label{eq:polar-left-module}
                  (J_L)_d=L\otimes\Sym^d\mathcal U
\end{equation}
as an actual subspace of the specified Cauchy component.  The ideal
has no lower-degree generators.  All hypotheses of
Theorem~\ref{thm:structural-coefficient-reconstruction} are now
verified, with \(m=d\) and \(\Lambda=\{(d)\}\), proving the
forward implication.  For the reverse implication the bundle map
\(T\mapsto gT\) transports both block maps, up to invertible
changes in their domain and target.  It therefore transports their
Fitting schemes.  Scalar choices of \(g\) do not alter the system.
\end{proof}

For a single equation \(L=\C f\), this yields an intrinsic
scheme-valued invariant of a degree-\(d\) hypersurface in
\(\Pj(V)\), up to projective equivalence.  For example \(n=3,d=4\)
gives all plane quartics, smooth or singular, using a degree-seven
Fitting ideal over \(\Gr(3,H)\).  Higher-dimensional systems give
Grassmannian families of equations rather than isolated examples.
The support is fixed throughout each such family, and the varying
coefficient system is recovered solely from its nonreduced structure.
This construction concerns arbitrary homogeneous systems.  The
original multiplication construction has the different natural
extension proved in Section~\ref{sec:natural-pencils}.
Nor does \eqref{eq:polar-system-torelli} assert an equivalence of
moduli stacks or a classification of all automorphisms of \(Z_L\);
it is the stated separation of abstract isomorphism classes.

\begin{example}[An explicit polar nilpotent layer]
\label{ex:polar-nilpotence-three}
For \(n=d=2\), \(f=x_1^2+x_2^2\), and
\(T=\left(\begin{smallmatrix}a&b\\c&e\end{smallmatrix}\right)\),
put \(\delta=ae-bc\), \(A=a^2+c^2\), \(B_0=ab+ce\),
and \(C=b^2+e^2\).  The polar ideal is
\(J=(A,2B_0,C)\), and \(I(Z_{\C f})=\delta J\).
The identity \(AC-B_0^2=\delta^2\) gives \(\delta^3\in I\).
But \(\delta\notin J_2=\langle A,B_0,C\rangle\), since the
monomials \(ae,bc\) do not occur in that span.  Cancellation in
the polynomial domain gives \(\delta^2\notin I\).
Since \(\sqrt I=(\delta)\), the nilradical therefore has
nilpotency index exactly three on this chart.  The example displays
scheme structure lost by the common reduced determinant support.
\end{example}
